\subsection{A dimension-150 central boson}
\label{sec:r160}

\begin{proposition}
The census ring $R_{7,2,160}$ has no categorification over $\C$.
\end{proposition}
\begin{proof}
Its basis is $1,g,v,X_1,\ldots,X_4$, with dimensions
$(1,1,2,6,6,6,6)$ and total dimension $150$. All basis objects are
self-dual, the adjoint subring is the whole ring, and
\[
g^2=1,\quad gv=v,\quad gX_i=X_i,\quad
v^2=1+g+v,\quad vX_i=2X_i.
\]
Suppose that $\mathcal C$ realizes this ring. It is integral and
pseudo-unitary \cite{ENO}. Its center is integral and modular. The
modular formal-codegree formula implies that every central simple
dimension $d$ divides $150$: the rational number $150^2/d^2$ is an
algebraic integer, so the rational number $150/d$ is an integer.

Write $F:\mathcal Z(\mathcal C)\to\mathcal C$ for the forgetful functor
and $I$ for its right adjoint. The formal codegrees are
$150,3,2,25,25,25,25$. Thus $I(1)$ has seven distinct constituents,
each of multiplicity one, with dimensions $1,50,75,6,6,6,6$
\cite[Theorem~2.13]{OstrikRank3}. Each restriction contains $1$ once.
The formula
$FI(W)=\bigoplus_{S\in\operatorname{Irr}(\mathcal C)}SWS^*$ gives
\[
\langle I(1),I(g)\rangle=5,\qquad
\langle I(1),I(v)\rangle=9,\qquad
\langle I(g),I(g)\rangle=7.
\]
For a dimension-six constituent $Y$ of $I(1)$, write
$F(Y)=1+b g+c v$; then $b+2c=5$, so $b$ is positive and odd.
The four such constituents use at least four of the total five copies
of $g$. Consequently all four have $b=1,c=2$.
If $b_{50},b_{75},c_{50},c_{75}$ are the corresponding coefficients
for the other two nontrivial constituents, the induction pairings give
$b_{50}+b_{75}=c_{50}+c_{75}=1$. Reduction of their dimensions modulo
six gives $1+b_{50}+2c_{50}\equiv50\pmod6$, hence
$(b_{50},c_{50},b_{75},c_{75})=(1,0,0,1)$.
It follows that $I(g)$ shares precisely five multiplicity-one
constituents with $I(1)$, of total dimension $74$.
Its squared norm is seven, and its dimension is $150$, so its other
two constituents are distinct, of multiplicity one, and have dimensions
summing to $76$. The only pair of positive divisors of $150$ with this
sum is $(1,75)$. A central invertible object $\widetilde g$ over $g$
is therefore forced.

The actual fusion subcategory $\mathcal D=\langle1,g,v\rangle$ has a
rank-one module category generated by any $X_i$. This gives a fiber
functor; semisimple Hopf reconstruction in dimension six identifies
$\mathcal D$ with $\operatorname{Rep}(S_3)$ \cite{EGpq}.
Its universal grading is trivial, so the half-braiding on $g$ is unique:
two choices differ by a tensor automorphism of the identity. Its
restriction is therefore the symmetric sign object, with self-braiding
$+1$. The universal grading of $\mathcal C$ is also trivial, so a
central invertible over $1$ is trivial. Hence $\widetilde g^2=1$, and
$\langle\widetilde g\rangle$ is a central
$\operatorname{Rep}(C_2)$.

De-equivariantization yields
$\mathcal C=\mathcal E^{C_2}$ with $\FPdim(\mathcal E)=75$.
The equivariantization orbit formula \cite{BurciuNatale} and the
sign action give $\mathcal E$ the type $1^3,3^8$: each sign-fixed
dimension-two or dimension-six source simple comes from a free orbit
of two objects of dimension one or three. For a dimension-three simple
$Z$ of $\mathcal E$, its invertible stabilizer has order one or three.
Order one would leave dimension eight in $ZZ^*$ to be supplied by
dimension-three objects, which is impossible. Thus $ZZ^*$ contains all
three invertibles and two dimension-three constituents.
The adjoint subcategory therefore has dimension $3+9k$ with $k>0$;
divisibility by $75$ forces it to equal $\mathcal E$.

By the two-prime theorem, $\mathcal E$ is solvable
\cite[Theorem~1.6]{ENOSolvable}. Its last prime step cannot be an
extension, since its universal grading is trivial; it is an
equivariantization by $C_3$, since its invertible group is $C_3$.
The resulting central $\operatorname{Rep}(C_3)$ has these invertibles
as its underlying objects. Their half-braidings are unique, so the
outer $C_2$ action preserves this central subcategory.
Equivariantizing it embeds a central symmetric category of dimension
six into $\mathcal Z(\mathcal C)$: its half-braidings with objects of
$\mathcal E$ descend to equivariant objects. Forgetting equivariance
gives a fiber functor, so this category is $\operatorname{Rep}(H)$ for
a group of order six. It has exactly two invertibles, because it
contains the canonical $\operatorname{Rep}(C_2)$ and embeds in
$\mathcal C$. Thus $H=S_3$, and its image is precisely
$\langle1,g,v\rangle$.
Direct de-equivariantization has dimension $25$, hence is pointed
\cite{ENO,EGO}. We obtain
\[
\mathcal C=(\operatorname{Vec}_A^\omega)^{S_3},\qquad |A|=25.
\]

For a nonidentity element of $A$ in an orbit of length $\ell<6$, a
simple lift has dimension $\ell\dim\pi$, where $\pi$ is an irreducible
projective representation of a stabilizer of order $6/\ell$.
It is at most $\sqrt{6\ell}<6$. Such a lift cannot belong to
$\langle1,g,v\rangle$, whose objects lie over the identity element of
$A$, and every other simple of $\mathcal C$ has dimension six.
Thus $S_3$ acts freely on $A\setminus\{1\}$.
If $A=C_{25}$, its automorphism group has order twenty and the
order-three subgroup acts trivially. If $A=C_5^2$, every involution
acting without nonzero fixed vectors is $-I$, since its minimal
polynomial divides $(t-1)(t+1)$. All three involutions would therefore
have the same image, making their products, the order-three elements,
act trivially. Both cases contradict freeness.
\end{proof}

The exact certificate checks the complete source tensor, the adjoint
subring, formal codegrees, the three induction pairings, and the small
integer and finite-group premises. It does not enumerate a full
induction matrix.
